% -----------------------------------------------------------------------------
% This file is automatically generated.
% Do not edit manually.
% Source: notebooks/support/regression-analysis/support.Rmd
% -----------------------------------------------------------------------------


\noindent We now turn to a model with multiple independent variables, for example, \ttblue{mod.1}, introduced in Section 5.4.3.

\par\addvspace{\topsep}
\begingroup
\fvset{listparameters={%
  \setlength{\topsep}{0pt}%
  \setlength{\partopsep}{0pt}%
  \setlength{\parsep}{0pt}%
  \setlength{\itemsep}{0pt}%
}}
\begin{Verbatim}[breaklines=true,formatcom=\color{ada_blue}]
(anova.mod.1 <- anova(ussteamco.mod.1))
\end{Verbatim}
\begin{Verbatim}[breaklines=true,formatcom=\color{ada_light_blue}]
Analysis of Variance Table

Response: revenue
           Df     Sum Sq    Mean Sq F value    Pr(>F)
production  1 3.7804e+14 3.7804e+14  96.361 3.561e-11 ***
coolDD      1 3.5110e+14 3.5110e+14  89.493 8.659e-11 ***
heatDD      1 9.8390e+13 9.8390e+13  25.079 1.942e-05 ***
Residuals  32 1.2554e+14 3.9232e+12
---
Signif. codes:  0 '***' 0.001 '**' 0.01 '*' 0.05 '.' 0.1 ' ' 1
\end{Verbatim}
\endgroup
\par\addvspace{\topsep}

\noindent Declaring the independent variables in a different order shows that the calculation of the sum of squares is in sequential order.
